\chapter{致\texorpdfstring{\quad}{}谢}

回首在华园的七年时光，心中满是感激。

首先谨向我的导师杨辰光教授致以最诚挚的谢意。杨老师性格儒雅随和，治学严谨，在我心中是一位高尚的学者。每当我在研究中遇到困惑，杨老师总是耐心地给予指导，其渊博的学识让我受益匪浅。更令我感念的是，杨老师待人接物始终宽厚温和，让我学会如何以善意和包容对待身边的人。感谢实验室的同门李沐、杰彬、宇博、正聪、深元、思睿、雅静和志毅，和你们一起谈天说地、一起攻克难关，这些经历让我的研究生生活变得丰富多彩。还要感谢从本科相识至今的舍友浩泽、潮海和方烁。我们在求职最艰难的时刻相互扶持，在苦闷时彼此鼓励，这份友谊是我求学生涯中珍贵的财富。

上海的经历仿佛是人生中一段小小的岔路，却也遇见了许多良师益友。感谢建锋，我们仨共患难的日子应是很难忘记了。感谢方正，你对我的帮助和倾囊相授的职场智慧，让我受益良多。感谢程冉，你的激情与鼓励让我明白，永远要心存梦想。感谢於辉、鸿飞和婉婷，你们让我学会如何成为一个务实而善良的人。还要感谢朱骞、泽辰、奕宇给予我的诸多帮助，以及查琰、赵山、柯宇、刘叙、豪杰、若澎、文群、所航、梓阳、艺严、泽宇，与你们相遇是一种难得的幸运。感谢一凡、黄维、林一、居垚，是你们让我在异乡减轻了许多孤单。

感谢从高中相识至今的好友广能、起源和楚翘。从高二到现在，转眼快是十年光景。从未想过，十年之后会一起在同济的球场上挥汗如雨，一起看很多演出，一起到肇庆乡村走访调研。时光飞逝，高中时代仿佛仍在昨天，我们的友谊始终如初。

最后，将最深的感谢献给我的父母。感谢你们一路走来的鼓励与支持，你们总是竭尽所能，让我能够心无旁骛地前行。你们的爱是我最坚实的后盾。

行文至此，才发现这路上承蒙太多人的帮助与善意，寥寥数语难以道尽心中感激。今天是个浪漫的日子，我希望对所有照亮过我的人说，我爱你们。

\begin{minipage}[t]{0.945\textwidth}%
	\begin{flushright}
		2026年5月20日\\
		于华南理工大学
		\par\end{flushright}
\end{minipage}

